\section{Introduction}

Scaling laws turn small experiments into forecasts of how loss changes with
model size, data, and compute.  Empirical power laws have been observed across
vision, language, speech, and generative modeling
\cite{hestness2017deep,rosenfeld2020constructive,kaplan2020scaling,hoffmann2022training},
while spectral theories explain many such curves through the progressive
resolution of weaker covariance or kernel modes
\cite{sollich1998learning,bordelon2020spectrum,canatar2021spectral,bahri2024explaining,lin2024scaling}.
These results support a useful planning picture: more data reveals increasingly
fine predictive structure, and the unresolved tail determines the learning
slope.

Most scaling-law theories formalize this picture with independent examples.
Yet many high-value datasets are contiguous trajectories: robot experience,
physical and environmental sensors, financial histories, user-interaction
logs, and replay buffers.  A common practical correction replaces the nominal
sample count $N$ by a scalar effective sample size $N_{\rm eff}$ computed from
a global mixing time or integrated autocorrelation.  This is appropriate when
all feature directions share one temporal kernel.  It can be misleading when
the directions that dominate late-stage learning are also the slowest to
refresh.

Consider two streams with exactly the same one-time covariance spectrum.  In
the first, every spectral direction persists for the same duration.  In the
second, high-index directions are individually weak but occur in long bursts.
A global average can assign the two streams the same autocorrelation time.
Nevertheless, late directions arrive much less often in the second stream, so
the number of observations needed to discover them grows with their index.
The relevant object is then not one effective sample count but a sequence of
mode-wise information rates.  This motivates our central question:
\begin{quote}
\emph{When does dependence change a scaling exponent, rather than only a
constant?}
\end{quote}

We answer this question in an exactly solvable persistent regression model.
Mode $j$ occupies a fraction $\lambda_j$ of raw time but lasts $\tau_j$ steps
whenever it appears.  Preserving the same one-time covariance forces a new
block to choose mode $j$ with probability
\begin{equation}
q_j=\frac{\lambda_j/\tau_j}{\sum_\ell\lambda_\ell/\tau_\ell}.
\label{eq:intro-q}
\end{equation}
Thus $\lambda$ is the \emph{marginal spectrum}, while $q$ is the
\emph{innovation spectrum}: it measures the arrival of statistically new
spectral information.  If $\lambda_j\asymp j^{-a}$ and
$\tau_j\asymp j^r$, then $q_j\asymp j^{-(a+r)}$.  After $B$ innovations, the
largest mode seen with constant probability is
$J_B\asymp B^{1/(a+r)}$.  If target energy satisfies
$s_j=\lambda_j\theta_j^2\asymp j^{-b}$, the unresolved tail is
\begin{equation}
\sum_{j>J_B}s_j\asymp B^{-(b-1)/(a+r)}.
\label{eq:intro-mechanism}
\end{equation}
This simple coverage mechanism changes the slope while leaving the marginal
prediction problem unchanged.

Our contributions are as follows.
\begin{enumerate}
\item We define \emph{spectral renewal regression} and derive an exact
finite-sample decomposition into model approximation, unseen-mode bias, and
label-noise variance.  Under power laws, all three terms have matching upper
and lower rates.
\item We prove an exact noiseless minimax equality over a signed source class.
No optimizer, number of epochs, or data-reuse strategy can infer the sign of a
predictive mode that never arrived.
\item We give a direct counterexample to scalar effective sample size.  Two
stationary streams have the same marginal covariance and the same trace
integrated autocorrelation time, yet have exponents $(b-1)/a$ and
$(b-1)/(a+r)$.  We then derive the sample- and compute-optimal consequences of
the mode-wise law.
\item We pass from a fixed number of innovations to exactly $N$ raw examples.
At a renewal boundary, the innovation exponent holds for every $r\ge0$
without a finite-variance dwell assumption.  Under stationary initialization,
a size-biased initial block can create a separate inspection-paradox phase.
\item Using the von Bahr--Esseen inequality, we transfer the noisy hard-target
rate under fractional, rather than second, dwell moments.  This covers
substantial infinite-variance regimes.
\item We prove invariance under known well-conditioned dense representations
and deliberately test the boundary of the theory.  Dense-dictionary
regression retains the predicted slopes, while a dense Gaussian
autoregressive negative control does not: autocorrelation alone is not enough
when every sample contains every mode.
\end{enumerate}

The distinction in the last item is important.  We do not claim that arbitrary
nonseparable time series universally replace $a$ by $a+r$.  The exponent
changes when persistence limits the arrival of new predictive spectral
information.  This narrower claim is both provable and operational: it suggests
estimating the persistence-adjusted spectrum $\lambda_j/\tau_j$, rather than
correcting all directions by one global mixing time.
